\documentclass{article}
\usepackage{amsmath, amssymb, amsthm}
\usepackage{hyperref}
\usepackage{geometry}
\geometry{margin=1in}

\title{Erd\H{o}s Problem \#404}
\date{Last updated: 2025-08-31}
\author{Source: erdosproblems.com}

\begin{document}
\maketitle

\noindent\textbf{Status:} OPEN \quad \textbf{No prize}

\noindent\textbf{Tags:} number theory, factorials

\medskip
\noindent\textbf{Problem Statement:}

For which integers $a\geq 1$ and primes $p$ is there a finite upper bound on those $k$ such that there are $a=a_1<\cdots<a_n$ with\[p^k \mid (a_1!+\cdots+a_n!)?\]If $f(a,p)$ is the greatest such $k$, how does this function behave?

Is there a prime $p$ and an infinite sequence $a_1<a_2<\cdots$ such that if $p^{m_k}$ is the highest power of $p$ dividing $\sum_{i\leq k}a_i!$ then $m_k\to \infty$?

\bigskip
\noindent\textbf{Remarks:}

See also [403]. Lin \cite{Li76} has shown that $f(2,2) \leq 254$.

\bigskip
\noindent\textbf{References:}

[Li76] Lin, S., On two problems of Erd\H{o}s concerning sums of distinct factorials. Bell Laboratories internal memorandum (1960).

\bigskip
\noindent\small{Source: \url{https://www.erdosproblems.com/404}}
\end{document}
